\section{Le Méta-Modèle Entity}\label{sub:ent}
Le méta-modèle Entity représente une couche \og métier \fg{}. Il permet notamment de modéliser la persistance des données :  

\begin{figure}[htb]
  \centering
  \includegraphics[scale=.4]{img/Entity.eps}
  \caption{Méta-modèle Entity}
  \label{fig:ent}
\end{figure}



\subsection{Gestion des entités avec \kwplay{}}
Le framework \kwplay{} utilise la notion de \textit{Model} qui s'accorde parfaitement avec le concept d'\kwentity{}. 

\subsection{Conception du modèle}
Nous avons employé la simple démarche suivante : pour chaque modèle on associe une entité. À ce stade, on en déduit facilement les attributs de chaque instance d'\verb+Entity+ et les instances \verb+Reference+ correspondant aux associations avec d'autres entités. \kwplay{} utilise des mécanismes d'annotation au sein de ses modèles. Ces annotations permettent notamment de donner des contraintes à des attributs. Nous avons facilement pris en compte ce concept en utilisant la classe \verb+Annotation+ mise à disposition dans les méta data d'un \verb+Attribut+.
\clearpage


% LocalWords:  Entity framework Model
